% Generated draft: requires mathematical and source review.
\documentclass[11pt]{article}
\usepackage[margin=25mm]{geometry}
\usepackage{fontspec}
\setmainfont{Latin Modern Roman}
\usepackage{amsmath,amssymb,mathtools,unicode-math}
\setmathfont{Latin Modern Math}
\usepackage{enumitem,needspace,xurl,hyperref}
\setlength{\parindent}{0pt}
\setlength{\parskip}{5pt}
\setlength{\emergencystretch}{3em}
\newcommand{\N}{\mathbb{N}}
\newcommand{\Z}{\mathbb{Z}}
\newcommand{\Q}{\mathbb{Q}}
\newcommand{\R}{\mathbb{R}}
\newcommand{\C}{\mathbb{C}}
\newcommand{\F}{\mathbb{F}}
\newcommand{\A}{\mathbb{A}}
\newcommand{\PP}{\mathbb{P}}
\newcommand{\E}{\mathbb{E}}
\newcommand{\eps}{\varepsilon}
\newcommand{\dd}{\,\mathrm{d}}
\begin{document}
\section*{UnsolvedMath ID: 30001936}
\textbf{Source problem code:} catalogue entry
\subsection{Definitions and mathematical statement}
A binary matroid is a matroid representable over the field \(\F_2\). A minor of a matroid is obtained by deletions and contractions. Write \(N\preccurlyeq M\) when \(M\) has a minor isomorphic to \(N\), and write \(M^*\) for the dual matroid. The affine geometry matroid \(\operatorname{AG}(3,2)\) is represented over \(\F_2\) by the eight columns \((1,x)^{\mathsf T}\), indexed by \(x\in\F_2^3\). The source also names a matroid \(T_{11}\) and its dual \(T_{11}^*\), but supplies no representation or other identifying definition of \(T_{11}\).

Let \(\mathcal B\) denote the class of binary matroids in the intended scope of the source, and let \(\operatorname{MFMC}_1(M)\) denote the assertion that \(M\) satisfies the max-flow min-cut relation intended in the source's use of ``\(1\)-flowing.'' The supplied formulation of Seymour's conjecture is
\[
\forall M\in\mathcal B,\qquad
\left(
\operatorname{AG}(3,2)\not\preccurlyeq M
\ \land\ T_{11}\not\preccurlyeq M
\ \land\ T_{11}^*\not\preccurlyeq M
\right)
\Longrightarrow \operatorname{MFMC}_1(M).
\]
The question is whether this implication holds for every such binary matroid.

\emph{Unresolved qualifications.} The record does not explicitly specify whether all matroids are finite. It also does not define the flow and cut objects, the admissible capacities, the integrality requirements, or the quantification over any distinguished element in the max-flow min-cut assertion. Consequently, \(\operatorname{MFMC}_1\) above is an explicitly unresolved source predicate, not a complete definition of the \(1\)-flowing property. Together with the missing identification of \(T_{11}\), these omissions prevent a fully self-contained precise formulation from the supplied evidence alone.
\subsection{Short English statement}
The supplied conjecture, attributed to Seymour, asks whether every binary matroid having none of \(\operatorname{AG}(3,2)\), \(T_{11}\), and \(T_{11}^*\) as a minor satisfies the intended \(1\)-flowing max-flow min-cut property. The record leaves that property and \(T_{11}\) undefined, so the displayed formulation retains these unresolved dependencies.
\subsection{Sources}
\begin{itemize}
\item Supplied local UnsolvedMath registry record, permanent ID \(30001936\), title ``Seymour One-Flowing Conjecture,'' source problem code \texttt{catalogue entry}, in the Oberwolfach Reports Open Problems collection. The conjecture is attributed to Seymour; the proposer field is empty, and the proposed year is recorded as 2011.
\item Supplied report citation: \emph{Combinatorial Optimization (2012)}, Oberwolfach Reports, DOI 10.4171/owr/2011/53, pp. 3003--3076. Original OWR report: \url{https://doi.org/10.4171/owr/2011/53}. The background describes these reports as records of mathematical workshops and their open-problem discussions. The registry labels the statement \texttt{corrected\_verified} and says the official report gives precisely the three listed excluded minors; this is supplied verification metadata, not an independent verification here.
\item Supplied literature source, labelled ``binary matroid templates'': \url{https://doi.org/10.1137/16M1077283}. No authors, full title, or precise supporting theorem are supplied.
\item Literature assessment dated \texttt{2026-08-21}: the supplied curation labels the problem ``partially solved'' and states that the report calls it Seymour's \(1\)-flowing conjecture, that later work characterizes highly connected cases, and that the excluded-minor conjecture is not fully resolved. This is supplied dated curation, not independently verified. The background expressly requests independent verification of the time-sensitive status and cited claims.
\end{itemize}
\end{document}
